\documentclass[12pt,a4paper,openany,oneside]{book}

% --- Paquetes ---
\usepackage[utf8]{inputenc}
\usepackage[spanish, es-noshorthands]{babel}
\usepackage{amsmath, amsfonts, amssymb}
\usepackage{graphicx}
\usepackage{geometry}
\usepackage{hyperref}
\usepackage{booktabs}
\usepackage{fancyhdr}
\usepackage{listings}
\usepackage{xcolor}
\usepackage{tikz}
\usepackage{pgfplots}
\usepackage{colortbl}
\usepackage{multirow}
\usepackage{hhline}
\usetikzlibrary{arrows.meta, positioning, shapes.geometric, calc, shapes, chains, scopes}
\pgfplotsset{compat=1.18}

\geometry{margin=2.5cm}

% --- Colores y estilos para código Python ---
\definecolor{codegreen}{rgb}{0,0.6,0}
\definecolor{codegray}{rgb}{0.5,0.5,0.5}
\definecolor{codepurple}{rgb}{0.58,0,0.82}
\definecolor{backcolour}{rgb}{0.95,0.95,0.92}

\lstset{
    language=Python,
    backgroundcolor=\color{backcolour},   
    commentstyle=\color{codegreen},
    keywordstyle=\color{blue},
    numberstyle=\tiny\color{codegray},
    stringstyle=\color{codepurple},
    basicstyle=\ttfamily\small,
    breaklines=true,
    frame=single,
    showstringspaces=false
}

% --- Estilo de cabecera y pie ---
\pagestyle{fancy}
\fancyhf{}
\renewcommand{\headrulewidth}{0.4pt}
\renewcommand{\footrulewidth}{0.4pt}
\fancyhead[L]{\small Carlos César Sánchez Coronel}
\fancyhead[R]{\small Sesión 10: Modelos Pre-entrenados (BERT, GPT)}
\fancyfoot[L]{\small Carlos César Sánchez Coronel}
\fancyfoot[C]{\thepage}

% --- Portada ---
\title{
    \textbf{Sesión 10} \\ 
    \Huge \textbf{MODELOS PRE-ENTRENADOS (BERT, GPT) Y FINE-TUNING} \\
    \Large La revolución del transfer learning en NLP
}
\author{
    \small Por \\ \vspace{0.2cm}
    \Large \textsc{Carlos César Sánchez Coronel}
}
\date{2026}

\begin{document}

\maketitle
\tableofcontents
\addcontentsline{toc}{chapter}{Índice general}

% ------------------------------------------------------------------
% CAPÍTULO 10: SESIÓN 10 - MODELOS PRE-ENTRENADOS
% ------------------------------------------------------------------
\chapter{Modelos Pre-entrenados (BERT, GPT) y Fine-tuning}

La arquitectura Transformer, presentada en la sesión anterior, no solo revolucionó el diseño de modelos, sino que también permitió el desarrollo de modelos pre-entrenados a gran escala. BERT (Bidirectional Encoder Representations from Transformers) y GPT (Generative Pre-trained Transformer) son los dos representantes más influyentes de esta nueva era. Esta sesión explora sus arquitecturas, objetivos de pre-entrenamiento, el proceso de fine-tuning para tareas específicas, y las técnicas modernas para adaptar eficientemente estos modelos a problemas concretos.

\section{Logro de la sesión}
Utilizar y fine-tunar modelos pre-entrenados como BERT y GPT para tareas específicas, comprendiendo sus arquitecturas, objetivos de entrenamiento y las diferencias fundamentales entre ellos.

\section{Introducción: el poder del pre-entrenamiento}
Antes de 2018, la mayoría de los modelos de NLP se entrenaban desde cero para cada tarea, lo que requería grandes cantidades de datos etiquetados. La introducción de modelos pre-entrenados como ELMo, ULMFiT, y especialmente BERT y GPT, cambió radicalmente este panorama. La idea es entrenar un modelo masivo en grandes corpus de texto no etiquetado (aprendiendo representaciones ricas del lenguaje) y luego adaptarlo (fine-tuning) a tareas específicas con relativamente pocos datos etiquetados. Esto se conoce como \textbf{transfer learning} y es el estándar actual en NLP.

\section{BERT (Bidirectional Encoder Representations from Transformers)}

BERT, propuesto por Devlin et al. (2018) , es un modelo basado en el encoder del Transformer. Su principal innovación es el entrenamiento bidireccional profundo, a diferencia de los modelos anteriores que solo leían texto de izquierda a derecha o combinaban direcciones de forma superficial.

\subsection{Arquitectura}
BERT utiliza solo la parte \textbf{encoder} del Transformer. Apila múltiples capas de encoder (12 para BERT-base, 24 para BERT-large). Cada capa contiene:
\begin{itemize}
    \item Multi-head self-attention (bidireccional, sin máscara causal).
    \item Feed-forward network.
    \item Layer normalization y conexiones residuales.
\end{itemize}

La entrada a BERT es una secuencia de tokens que puede ser una o dos oraciones (separadas por un token especial `[SEP]`). Se añade un token especial `[CLS]` al inicio, cuya representación final se utiliza para tareas de clasificación.

\begin{center}
\begin{tikzpicture}[
    layer/.style={rectangle, draw, minimum width=3cm, minimum height=0.8cm, fill=blue!10},
    arrow/.style={-latex, thick}
]
\node[layer] (embed) at (0,0) {Embedding + Posición};
\node[layer] (enc1) at (0,1.5) {Encoder 1};
\node[layer] (enc2) at (0,3) {Encoder 2};
\node[layer] (encN) at (0,4.5) {Encoder 12};
\node[layer] (cls) at (0,6) {Representación [CLS]};

\draw[arrow] (embed) -- (enc1);
\draw[arrow] (enc1) -- (enc2);
\draw[arrow] (enc2) -- (encN);
\draw[arrow] (encN) -- (cls);
\end{tikzpicture}
\caption{Arquitectura simplificada de BERT.}
\end{center}

\subsection{Tokenización: WordPiece}
BERT utiliza el tokenizador \textbf{WordPiece}, que segmenta palabras en subpalabras basándose en la frecuencia de co-ocurrencia. Por ejemplo, "playing" podría dividirse en "play" + "##ing". Esto permite manejar palabras desconocidas y reduce el tamaño del vocabulario.

\textbf{Ejemplo}:
Entrada: "I love NLP"
Tokens: ["[CLS]", "I", "love", "NL", "##P", "[SEP]"]

\subsection{Objetivos de pre-entrenamiento}
BERT se entrena con dos objetivos simultáneos:

\subsubsection{Masked Language Model (MLM)}
Se enmascara aleatoriamente el 15\% de los tokens de entrada (reemplazándolos con `[MASK]`). El modelo debe predecir los tokens originales basándose en el contexto bidireccional. Esto fuerza al modelo a aprender representaciones contextuales profundas.

\textbf{Ejemplo}:
Frase original: "El gato persigue al ratón"
Frase enmascarada: "El gato [MASK] al ratón"
Objetivo: predecir "persigue".

\subsubsection{Next Sentence Prediction (NSP)}
Se entrena al modelo para predecir si dos oraciones son consecutivas en el texto original. Esto ayuda a tareas que requieren comprensión de relaciones entre oraciones (como pregunta-respuesta o inferencia). Se toman pares (A, B) donde el 50\% de las veces B es la oración siguiente real, y el otro 50\% es una oración aleatoria.

\textbf{Ejemplo}:
- A: "El gato persigue al ratón."
- B (siguiente): "El ratón se esconde."
- Etiqueta: "IsNext"
- B (aleatoria): "El coche es rojo."
- Etiqueta: "NotNext"

\subsection{Representaciones de salida}
Para cada token de entrada, BERT produce un vector contextualizado (última capa). La representación del token `[CLS]` se usa como representación de toda la secuencia para tareas de clasificación. Para tareas de etiquetado (NER, POS), se usan las representaciones de cada token.

\section{GPT (Generative Pre-trained Transformer)}

GPT, desarrollado por Radford et al. (2018) y sus versiones posteriores (GPT-2, GPT-3, GPT-4), es un modelo basado en el \textbf{decoder} del Transformer. Su objetivo es modelar la probabilidad de una secuencia de forma autorregresiva (predictiva hacia adelante).

\subsection{Arquitectura}
GPT utiliza solo la parte \textbf{decoder} del Transformer, pero con una modificación: en lugar de usar atención cruzada (pues no hay encoder), usa self-attention con máscara causal (solo puede mirar hacia atrás). Apila múltiples capas de decoder.

\begin{center}
\begin{tikzpicture}[
    layer/.style={rectangle, draw, minimum width=3cm, minimum height=0.8cm, fill=green!10},
    arrow/.style={-latex, thick}
]
\node[layer] (embed) at (0,0) {Embedding + Posición};
\node[layer] (dec1) at (0,1.5) {Decoder 1};
\node[layer] (dec2) at (0,3) {Decoder 2};
\node[layer] (decN) at (0,4.5) {Decoder 12};
\node[layer] (lm) at (0,6) {Capa de salida (LM head)};

\draw[arrow] (embed) -- (dec1);
\draw[arrow] (dec1) -- (dec2);
\draw[arrow] (dec2) -- (decN);
\draw[arrow] (decN) -- (lm);
\end{tikzpicture}
\caption{Arquitectura simplificada de GPT (decoder-only).}
\end{center}

\subsection{Tokenización: BPE (Byte-Pair Encoding)}
GPT utiliza \textbf{BPE}, similar al visto en sesiones anteriores. BPE fusiona los pares de bytes (o caracteres) más frecuentes hasta alcanzar el tamaño de vocabulario deseado. Es especialmente adecuado para modelar lenguaje porque maneja bien palabras raras y morfología.

\subsection{Objetivo de pre-entrenamiento}
GPT se entrena con el objetivo de \textbf{modelado de lenguaje autorregresivo}: dada una secuencia de tokens \(x_1, x_2, ..., x_T\), maximiza la probabilidad de cada token dado los anteriores:
\[
\mathcal{L} = \sum_{t=1}^T \log P(x_t | x_{<t})
\]
No hay enmascaramiento ni objetivo de predicción de siguiente oración. El modelo aprende a predecir la siguiente palabra de forma causal.

\subsection{Generación}
GPT es inherentemente generativo. Dado un prompt inicial, puede generar texto continuando la secuencia, muestreando de la distribución de probabilidad en cada paso. Esto lo hace ideal para tareas como:
\begin{itemize}
    \item Generación de texto creativo.
    \item Chatbots y diálogo.
    \item Completar código.
    \item Traducción (con prompting adecuado).
\end{itemize}

\section{Comparación BERT vs. GPT}

\begin{table}[h]
\centering
\begin{tabular}{|l|l|l|}
\hline
\textbf{Característica} & \textbf{BERT} & \textbf{GPT} \\
\hline
Arquitectura & Encoder-only & Decoder-only \\
Direccionalidad & Bidireccional (self-attention sin máscara) & Unidireccional (causal, con máscara) \\
Objetivo de entrenamiento & MLM + NSP & Modelado autorregresivo \\
Tokenización & WordPiece & BPE \\
Uso principal & Comprensión, clasificación, extracción & Generación, diálogo, completado \\
Fine-tuning & Requiere cabezas específicas (clasificación, etc.) & Puede fine-tunearse para tareas específicas o usarse con prompting \\
Ejemplos & BERT, RoBERTa, ALBERT & GPT-2, GPT-3, GPT-4 \\
\hline
\end{tabular}
\caption{Comparación entre BERT y GPT.}
\end{table}

\section{Fine-tuning vs. Feature-based Approaches}

\subsection{Feature-based}
En este enfoque, se utilizan los embeddings pre-entrenados como características fijas para un modelo más simple (ej. SVM, MLP). Por ejemplo, se pueden extraer las representaciones de BERT para cada oración y luego entrenar un clasificador lineal. Ventaja: más rápido y menos propenso a sobreajuste con datos muy pequeños. Desventaja: no se ajusta la representación a la tarea.

\subsection{Fine-tuning}
Aquí se añade una nueva capa (cabeza) específica para la tarea y se entrenan todos los parámetros del modelo (o algunos) conjuntamente. El gradiente fluye hacia atrás ajustando las representaciones pre-entrenadas a la tarea. Esto suele dar mejores resultados si se tienen suficientes datos.

\begin{center}
\begin{tikzpicture}[
    block/.style={rectangle, draw, minimum width=2.5cm, minimum height=1cm},
    arrow/.style={-latex, thick}
]
\node[block] (bert) at (0,0) {BERT pre-entrenado};
\node[block] (head) at (0,2) {Cabeza de clasificación};
\node[block] (loss) at (0,4) {Función de pérdida};
\node[block] (labels) at (2,4) {Etiquetas};

\draw[arrow] (bert) -- (head);
\draw[arrow] (head) -- (loss);
\draw[arrow] (labels) -- (loss);
\draw[arrow] (loss) -- (bert) node[midway, left] {gradientes};
\end{tikzpicture}
\caption{Fine-tuning: el gradiente ajusta los pesos de BERT y la cabeza.}
\end{center}

\subsection{Tipos de fine-tuning}
\begin{itemize}
    \item \textbf{Completo}: se actualizan todos los parámetros del modelo. Requiere más memoria y datos.
    \item \textbf{Con capas congeladas}: se congelan las primeras capas (que aprenden características generales) y solo se ajustan las últimas capas y la cabeza. Útil con datos pequeños.
    \item \textbf{Adaptadores (Adapters)}: se insertan pequeñas capas entrenables entre las capas pre-entrenadas. Solo se entrenan los adaptadores y la cabeza.
    \item \textbf{LoRA (Low-Rank Adaptation)}: se añaden matrices de bajo rango a las matrices de pesos, entrenando solo esas. Es extremadamente eficiente en parámetros.
\end{itemize}

\section{Técnicas avanzadas de fine-tuning}

\subsection{LoRA (Low-Rank Adaptation)}
LoRA, propuesto por Hu et al. (2021), congela los pesos pre-entrenados y añade matrices de descomposición de bajo rango a las matrices de atención y feed-forward. Si una matriz de pesos \(\mathbf{W}_0 \in \mathbb{R}^{d \times k}\) está congelada, LoRA aprende \(\mathbf{B} \in \mathbb{R}^{d \times r}\) y \(\mathbf{A} \in \mathbb{R}^{r \times k}\) con \(r \ll \min(d,k)\), de modo que la nueva matriz es \(\mathbf{W} = \mathbf{W}_0 + \mathbf{B}\mathbf{A}\). Solo se entrenan \(\mathbf{B}\) y \(\mathbf{A}\), reduciendo drásticamente el número de parámetros entrenables.

\textbf{Ventajas}:
- Permite fine-tunear modelos enormes (como GPT-3) con recursos limitados.
- Facilita el intercambio de adaptadores para diferentes tareas sin duplicar el modelo base.

\subsection{Adapters}
Se insertan pequeños módulos (normalmente dos capas feed-forward con una activación no lineal) entre las capas del transformer. Durante el fine-tuning, solo se entrenan estos adaptadores y la cabeza de tarea. Es similar a LoRA pero con una arquitectura diferente.

\subsection{Prompt Tuning}
En lugar de fine-tunear los pesos del modelo, se aprenden "prompts" continuos (vectores de soft prompts) que se anteponen a la entrada. El modelo permanece congelado. Es muy eficiente y funciona bien para modelos grandes como GPT-3.

\section{Tokenizadores de subpalabras en profundidad}

\subsection{WordPiece (usado en BERT)}
Algoritmo similar a BPE pero con una función de puntuación basada en la probabilidad del lenguaje:
\[
\text{score}(a,b) = \frac{\text{freq}(ab)}{\text{freq}(a) \cdot \text{freq}(b)}
\]
Se fusiona el par que maximiza esta puntuación. Esto tiende a crear tokens más lingüísticamente significativos.

\subsection{BPE (usado en GPT)}
Como se vio en sesiones anteriores, BPE fusiona iterativamente los pares de caracteres (o bytes) más frecuentes. Es más simple y muy efectivo.

\subsection{SentencePiece}
Implementación que no requiere pre-tokenización (trata el texto como secuencia de caracteres Unicode). Se usa en modelos como T5 y XLNet. Puede usar tanto BPE como unigram language model.

\section{Métricas de evaluación}

\subsection{Para clasificación}
\begin{itemize}
    \item \textbf{Accuracy}: porcentaje de predicciones correctas.
    \item \textbf{F1-score}: media armónica de precisión y recall, útil para clases desbalanceadas.
    \item \textbf{AUC-ROC}: área bajo la curva ROC, mide capacidad discriminativa.
\end{itemize}

\subsection{Para generación}
\begin{itemize}
    \item \textbf{Perplejidad}: exponencial de la entropía cruzada, mide qué tan "sorprendido" está el modelo. Baja perplejidad = mejor modelo.
    \item \textbf{BLEU}: para traducción, como se vio.
    \item \textbf{ROUGE}: para summarization, mide solapamiento de n-gramas entre generado y referencia.
\end{itemize}

\section{Aplicaciones prácticas}

\subsection{Clasificación de sentimiento con BERT}
Fine-tuning de BERT para clasificar reseñas como positivas o negativas. La cabeza es una capa lineal sobre la representación `[CLS]`.

\subsection{Reconocimiento de entidades nombradas (NER)}
Se añade una capa de clasificación por token sobre las representaciones de salida de BERT. Cada token se etiqueta como PERSON, LOC, ORG, etc.

\subsection{Respuesta a preguntas (QA)}
BERT se fine-tunea para predecir el inicio y fin de la respuesta en un párrafo dado. Se añaden dos cabezas lineales que predicen probabilidades de inicio y fin.

\subsection{Generación de texto con GPT}
Dado un prompt, GPT genera texto autoregresivamente. Con fine-tuning en un dominio específico, puede generar correos, artículos, código, etc.

\section{Ejemplo paso a paso: Fine-tuning de BERT para clasificación con Hugging Face}

\subsection{Paso 1: Instalación y carga de librerías}
\begin{lstlisting}[language=Python]
!pip install transformers datasets evaluate

from transformers import AutoTokenizer, AutoModelForSequenceClassification, Trainer, TrainingArguments
from datasets import load_dataset
import evaluate
import numpy as np
\end{lstlisting}

\subsection{Paso 2: Cargar datos}
Usaremos el dataset de IMDb para clasificación de sentimiento.
\begin{lstlisting}[language=Python]
dataset = load_dataset("imdb")
print(dataset)
\end{lstlisting}

\subsection{Paso 3: Tokenización}
\begin{lstlisting}[language=Python]
tokenizer = AutoTokenizer.from_pretrained("bert-base-uncased")

def tokenize_function(examples):
    return tokenizer(examples["text"], padding="max_length", truncation=True, max_length=512)

tokenized_datasets = dataset.map(tokenize_function, batched=True)
\end{lstlisting}

\subsection{Paso 4: Cargar modelo pre-entrenado}
\begin{lstlisting}[language=Python]
model = AutoModelForSequenceClassification.from_pretrained("bert-base-uncased", num_labels=2)
\end{lstlisting}

\subsection{Paso 5: Definir métrica}
\begin{lstlisting}[language=Python]
metric = evaluate.load("accuracy")

def compute_metrics(eval_pred):
    logits, labels = eval_pred
    predictions = np.argmax(logits, axis=-1)
    return metric.compute(predictions=predictions, references=labels)
\end{lstlisting}

\subsection{Paso 6: Configurar entrenamiento}
\begin{lstlisting}[language=Python]
training_args = TrainingArguments(
    output_dir="./results",
    evaluation_strategy="epoch",
    save_strategy="epoch",
    learning_rate=2e-5,
    per_device_train_batch_size=8,
    per_device_eval_batch_size=8,
    num_train_epochs=3,
    weight_decay=0.01,
    load_best_model_at_end=True,
    metric_for_best_model="accuracy",
)

trainer = Trainer(
    model=model,
    args=training_args,
    train_dataset=tokenized_datasets["train"].select(range(2000)),  # reducido para ejemplo
    eval_dataset=tokenized_datasets["test"].select(range(500)),
    compute_metrics=compute_metrics,
)
\end{lstlisting}

\subsection{Paso 7: Entrenar}
\begin{lstlisting}[language=Python]
trainer.train()
\end{lstlisting}

\subsection{Paso 8: Evaluar}
\begin{lstlisting}[language=Python]
trainer.evaluate()
\end{lstlisting}

\section{Preparación para entrevistas}

\subsection{Preguntas típicas}
\begin{itemize}
    \item \textbf{"¿Cuál es la diferencia entre BERT y GPT?"}
    BERT es encoder-only, bidireccional, entrenado con MLM y NSP; GPT es decoder-only, unidireccional, entrenado con modelado autorregresivo. BERT es mejor para comprensión, GPT para generación.
    
    \item \textbf{"¿Qué es el fine-tuning y cómo se diferencia de feature extraction?"}
    Fine-tuning ajusta todos los pesos del modelo a la tarea; feature extraction usa representaciones fijas. Fine-tuning suele dar mejores resultados si hay suficientes datos.
    
    \item \textbf{"¿Qué es LoRA y por qué es útil?"}
    Técnica de fine-tuning eficiente que añade matrices de bajo rango a las matrices de pesos. Reduce drásticamente el número de parámetros entrenables, permitiendo adaptar modelos grandes con menos recursos.
    
    \item \textbf{"¿Por qué BERT necesita positional encoding?"}
    Porque la atención es invariante a la posición; sin positional encoding, el modelo no sabría el orden de las palabras.
    
    \item \textbf{"¿Qué son MLM y NSP?"}
    MLM (Masked Language Model) predice palabras enmascaradas; NSP (Next Sentence Prediction) predice si dos oraciones son consecutivas. Son los objetivos de pre-entrenamiento de BERT.
    
    \item \textbf{"¿Cómo maneja BERT palabras desconocidas?"}
    Usa tokenizador WordPiece que descompone palabras en subpalabras; tokens desconocidos se dividen en subpalabras conocidas o se usa [UNK] si no hay.
\end{itemize}

\subsection{Preguntas avanzadas}
\begin{itemize}
    \item \textbf{"¿Qué es el problema de la máscara en MLM y cómo se soluciona en la práctica?"}
    Solo el 15% de los tokens se enmascaran; de ellos, el 80% se reemplazan con [MASK], 10% con token aleatorio, 10% se dejan igual. Esto evita la discrepancia entre pre-entrenamiento y fine-tuning (donde no hay [MASK]).
    
    \item \textbf{"¿Cómo funciona el attention mask en transformers?"}
    Máscara de atención indica qué tokens deben ignorarse (padding) o, en decodificadores, qué tokens futuros deben ocultarse (máscara causal).
    
    \item \textbf{"¿Qué son los embeddings de segmento en BERT?"}
    Se añaden embeddings que indican si un token pertenece a la primera o segunda oración (para tareas con pares).
\end{itemize}

\section{Conexión con el resto del curso}
\begin{itemize}
    \item \textbf{Sesión 8 (Seq2Seq y Atención)}: Base de los transformers.
    \item \textbf{Sesión 9 (Transformers)}: La arquitectura fundamental.
    \item \textbf{Sesiones anteriores}: Tokenización, embeddings, fine-tuning con modelos más simples.
\end{itemize}

\section{Resumen}
\begin{itemize}
    \item \textbf{BERT} es un modelo encoder-only bidireccional, pre-entrenado con MLM y NSP. Excelente para tareas de comprensión.
    \item \textbf{GPT} es un modelo decoder-only unidireccional, pre-entrenado con modelado autorregresivo. Ideal para generación.
    \item El \textbf{fine-tuning} adapta modelos pre-entrenados a tareas específicas, superando ampliamente a los enfoques clásicos.
    \item Técnicas como \textbf{LoRA} y \textbf{adaptadores} permiten fine-tuning eficiente con pocos recursos.
    \item Los \textbf{tokenizadores de subpalabras} (WordPiece, BPE) son esenciales para manejar vocabulario abierto.
    \item Estos modelos son la base de la mayoría de las aplicaciones modernas de NLP.
\end{itemize}

\end{document}